\section{Two exact conductors for the remaining small-$\eta$ chart}
These identities resolve specified higher-angle terms; they do not
constitute an admission of C14.3b. Keep
\[
 c^2=1/(\eta\zeta),\quad k^2=x,\quad y=x/c^2,\quad
 A=uv,\quad B=ckv^2/u,\quad C=k^2/c,\quad G=kv,
\]
\[
 d=B-G,\quad J=A-C,\quad F=V-y,\quad
 p=ckv/u,\quad R_c=pE+d.
\]
The elementary conductor is
\[
                         ud=k(cF-J).                 \tag{SC1}
\]
Consequently, for the pure-source coefficient $e_4=a_{UVV}$,
\[
 e_4xvEF=e_4(k/c)uR_cF-e_4xF^2+e_4(x/c)JF.           \tag{SC2}
\]
This is useful after the full angular completion. The exact scalar
remainder obeys
\[
 (c_F+e_4x)F^2+e_4vR_DF
       =(c_F+e_4(U+x))F^2-e_4yuH_uF.
\]
Substituting
$H_u=Q/\sqrt q-a(v/u)E/\eta+a\xi$ and using SC2 gives
\[
\begin{split}
 &(c_F+e_4U)F^2-e_4yuQF/\sqrt q-e_4ya u\xi F\\
 &\qquad+e_4(k/c)uR_cF+e_4(x/c)JF.
\end{split}                                             \tag{SC3}
\]
Here $a=1/\zeta$. On the source chart with $\eta,x$ infinitesimal
and $\zeta$ a positive unit,
\[
 \eta\lambda\sim\tfrac12,\quad
 \eta^2c_F\sim\tfrac12,\quad \eta^2e_4\sim\tfrac12,
 \quad q/x^2\text{ is a positive bounded unit}.
\]
Thus the normalized coefficients of the last two terms on
$(\sqrt\lambda R_c,\sqrt\lambda J,\sqrt{c_F}F)$ are
$O(\sqrt x\,u)$ and $O(x)$. The $Q,F$ coefficient is $O(u)$;
the $\xi,F$ coefficient is $O(xu)$, including normalization by
$\sqrt{B_x}$ since $B_x\ge1/2$. Retain $e_4UF^2$ as an exact
square. Finally $\sqrt\lambda R_c=R+\sqrt\lambda kE$ is
bounded on the coupled root $R$ and the weighted hard direction
$E/\sqrt\eta$. No bound on the bare $d$ is used.

For the complete divergent $D$ event, write the source block in
terms of $P=(1-r)E/\eta-\nu$, $Q_t=rE/\eta+\nu$, and
$b=Ar/\eta-(m\theta+\alpha+1/2)=O(x)$. Put
$T_D=R_D-uv\xi/h$. Then the complete original $D_p=D+1/2$ event is
$D_p\{P-hT_D-(j/\eta)vZ\}^2$. Its centered quadratic packet is
$D_p(P-hT_D)^2$; the contact-linear and scalar terms involving
$(j/\eta)vZ$ must remain in the exact contact balances. Replacing
$R_D$ by $T_D$ in the centered cross term $2bQ_tR_D$ leaves $2(b/h)uvQ_t\xi$. The exact
allocation is
\[
 uvQ_t=\frac{rU}{\eta ck}R_c-\frac{ru}{\eta}F
                    +\frac{ru}{\eta c}J+uv\nu.       \tag{SC4}
\]
After multiplication by $b$, the coefficients on
$\sqrt\lambda R_c,\sqrt{c_F}F,\sqrt\lambda J,\nu/\sqrt\eta$
are respectively $O(\sqrt x\,U),O(xu),O(xu),O(xuv\sqrt\eta)$.
Thus this particular full-$D$ cross term is harmless without comparing
$uv$ to $\sqrt\eta$.

\paragraph{Exact full angular cancellation.}
Put $D_0=u-cv$ and write the full angular root as $\mathcal Q$.
There are exact identities
\[
 R_c=kw(cE-D_0),\qquad
 \mathcal Q/\sqrt q=-D_0+(c/k)R_c-a\xi.
\]
Thus the complete principal $\xi$ row cancels to
\[
 -2\eta\lambda ckR_c\xi+2\eta\lambda c(R_D-uJ)\xi
 =-2\eta\lambda x\mathcal Q\xi/\sqrt q
                     -2\eta\lambda xa\xi^2.          \tag{SC5}
\]
Together with the singular rank's $+\mu\xi^2$, the last term leaves
$-\mu\xi^2$, since $\mu=\eta\lambda ax$; this is infinitesimal
relative to $B_x\xi^2$. The $\mathcal Q,\sqrt{B_x}\xi$ block
has strictly positive residue. Indeed, the exact source limits give
\[
 B_xq/x^2\sim
 \frac{B_x\zeta(1+\zeta)}2+\frac{(1+\zeta)^2}4
 >\frac14\sim(\eta\lambda)^2,
\]
where for unbounded $B_x$ the normalized cross tends to zero. For
bounded $B_x\ge1/2$, the difference has the explicit positive lower
bound $\zeta(3+2\zeta)/4$. No inequality
$B_x>1/(2\zeta)$ is needed.

\paragraph{The complete radial group, and the remaining dependency.}
Define $\overline F=F+vE$ and
$H_{\rm full}=\mathcal Q/\sqrt q+a\xi=H_u+c^2wE$.
Then exactly
\[
 uR_c=k(c\overline F-J),\qquad
 U-x=c^2\overline F-uH_{\rm full}.                    \tag{SC6}
\]
The entire $c_F$ group is $c_F\overline F^2$. The entire
$e_4$ group, including both hard and first-normal terms, has the
following exact simplification:
\[
\begin{split}
 &e_4\{U\overline F^2+y(U-x)(\overline F+vE)\}\\
 &=e_4(U+x)\overline F^2
   -e_4yu(\mathcal Q/\sqrt q+a\xi)\overline F
   +e_4yER_D.
\end{split}                                             \tag{SC7}
\]
The normalized $\mathcal Q,\overline F$ and
$\xi,\overline F$ coefficients are $O(u)$ and $O(xu)$.
The last term is an $O(x)$ coefficient perturbation of the weighted
source block because $e_4y\eta=O(x)$. It should be incorporated
before applying the full $D_p$ state shift.

There remains a specific root dependency: SC4 allocates the shifted
$uvQ_t\xi$ term on the \emph{bare} $F$ root. Completing the full
radial square retains $\overline F$, not an independent strict $F$.
Replacing $F=\overline F-vE$ in SC4 merely yields a tautology by
SC6. Therefore SC4 and SC7 cannot simply be joined while counting
both radial pivots. A possible completion must center the $m$ event jointly with the
other packets. The whole original $mR_m^2$ cannot be extracted as a
standalone positive square: at actual source contact
$R_m=\zeta M uZ\ne0$, although the parent vanishes. The relevant
exact centered identity is
\[
 R_m-\zeta M uZ
 =\frac{\mathcal Q}{a\sqrt q}
        +\frac r\eta u\widetilde F,\qquad
 \widetilde F=\frac\eta r vY-yZ,\quad
 M=1+xr-a\theta.                                      \tag{SC8}
\]
It preserves contact and indicates where the $uvQ_t$ control could
enter a simultaneous completion; the favorable source quotient
$b/\sqrt m=O(\sqrt x)$ alone does not supply that completion.
The contributions $2mr/\eta^2$ and $mr^2/\eta^2$ already counted
in the radial groups must also be reconciled exactly. Such a
contact-centered joint reassembly has not been proved here.
The displayed identities resolve the principal angular and full radial
obstructions but do not yet admit C14.3b.
